\scanpage{120}

\[
\left|\begin{aligned}
&n'\ne1,2,\quad\text{d'ordre }2p\text{ si }n'=2\quad(\text{donc }p\ne2\text{ !})\\
&\text{d'ordre }p\text{ si }n'=1.
\end{aligned}\right.
\]

\uline{Dém}. Posons en effet $\alpha=\unsure{\zeta|\int|\xi}+\zeta^{-1}$, où $\zeta$ appartient à un surcorps de $k$. La relation $\Psi_n(\alpha)=0$ signifie $\Phi_n(\zeta)=0$, et équivaut aussi, grâce à la proposition, à $\Phi_{n'}(\zeta)=0$, \deleted{Si $n'\geq3$,} \inserted{i.e. $\zeta$ est} \deleted{Celà signifie que $\Psi_{n'}(\alpha)=0$ donc que} \inserted{une racine de 1 d'ordre $n'$. Si} $n'\geq3$, celà signifie que l'ordre du polygône géométrique régulier d'invariant $\alpha$ est égal à $n'$, Si $n'=2$, que son ordre est $2p$, si $n'=1$ que son ordre est $p$, \hfill cqfd.

% archive sheet number: 119
